% ══════════════════════════════════════════════════════════════════
% Section 3 — Neural Network Architectures
% ══════════════════════════════════════════════════════════════════
\section{Neural network architectures}
\label{sec:architectures}

All architectures learn a scalar strain energy function
$\hat{\sef}: \R^d \to \R$, where $d = 3$ for isotropic materials
($\bar{I}_1, \bar{I}_2, J$) and $d = 5$ for single-fiber anisotropic
materials ($\bar{I}_1, \bar{I}_2, J, I_4, I_5$).  Stress-like quantities are
obtained by automatic differentiation:
$\partial\hat{\sef}/\partial I_\alpha$.

\subsection{Multilayer perceptron (MLP)}
\label{sec:architectures:mlp}

The baseline architecture is a standard feedforward network:
\begin{equation}
  \tens{z}^{(0)} = \tens{x}, \qquad
  \tens{z}^{(\ell)} = \phi\bigl(\mat{W}^{(\ell)}\tens{z}^{(\ell-1)}
    + \tens{b}^{(\ell)}\bigr), \qquad
  \hat{\sef} = \mat{w}^{(L)}\tens{z}^{(L-1)} + b^{(L)},
  \label{eq:mlp}
\end{equation}
where $\phi$ is the softplus activation $\phi(x) = \ln(1 + e^x)$, chosen for
its smoothness (required for second-derivative computations in the tangent
operator).  Unless stated otherwise, we use two hidden layers of width~64.

The MLP imposes no convexity or symmetry constraints and serves as a universal
approximation baseline.

\subsection{Input-convex neural network (ICNN)}
\label{sec:architectures:icnn}

Following Amos et al.~\cite{% TODO: Amos2017
}, the ICNN guarantees that $\hat{\sef}$ is convex in $\tens{x}$ by
construction:
\begin{align}
  \tens{z}^{(1)} &= \phi\bigl(\mat{W}^{(1)}_x \tens{x} + \tens{b}^{(1)}\bigr), \nonumber \\
  \tens{z}^{(\ell)} &= \phi\bigl(
    \underbrace{\mat{\tilde{W}}^{(\ell)}_z}_{\geq 0}\tens{z}^{(\ell-1)}
    + \mat{W}^{(\ell)}_x \tens{x}
    + \tens{b}^{(\ell)}\bigr), \qquad \ell = 2, \ldots, L, \nonumber \\
  \hat{\sef} &= \underbrace{\tilde{\mat{w}}^{(L)}_z}_{\geq 0} \tens{z}^{(L-1)}
    + \mat{w}^{(L)}_x \tens{x} + b^{(L)}.
  \label{eq:icnn}
\end{align}
Non-negativity of the z-path weights is enforced via a softplus
reparameterization: $\mat{\tilde{W}} = \ln(1 + e^{\mat{W}_\text{raw}})$.
Skip connections from the input $\tens{x}$ to every hidden layer ensure
expressive power despite the convexity constraint.

Convexity of $\hat{\sef}$ with respect to the invariants is a necessary (but
not sufficient) condition for material stability.  The ICNN output is always
scalar ($\hat{\sef} \in \R$).

\subsection{Polyconvex ICNN}
\label{sec:architectures:polyconvex}

Polyconvexity --- convexity of $\sef$ jointly in $(\Fmat, \text{cof}\,\Fmat,
\det\Fmat)$ --- is a stronger thermodynamic requirement that implies
quasi-convexity and ellipticity \cite{% TODO: Ball1977, Schroder2003
}.  Following Klein et al.~\cite{% TODO: Klein2022
}, we decompose the invariant space into groups and assign a separate ICNN
branch to each:
\begin{equation}
  \hat{\sef}(\tens{x}) = \sum_{g=1}^{G}
    \hat{\sef}_g\bigl(\tens{x}_{g}\bigr),
  \label{eq:polyconvex}
\end{equation}
where each $\hat{\sef}_g$ is an ICNN that is convex in its respective
invariant group $\tens{x}_g$.  For isotropic materials, we use three groups:
$\{I_1\}$, $\{I_2\}$, $\{J\}$.  Each branch is independently convex, and
their sum preserves convexity.

The total energy inherits polyconvexity provided the invariant groups are
chosen to correspond to convex functions of $(\Fmat, \text{cof}\,\Fmat,
\det\Fmat)$, as shown in \Cref{tab:polyconvex_groups}.

\begin{table}[htbp]
\centering
\caption{Polyconvex ICNN invariant groups.}
\label{tab:polyconvex_groups}
\begin{tabular}{@{}lll@{}}
\toprule
\textbf{Material type} & \textbf{Groups} & \textbf{Convexity argument} \\
\midrule
Isotropic   & $\{\bar{I}_1\}$, $\{\bar{I}_2\}$, $\{J\}$
  & Each is a convex function of $\Fmat$, $\text{cof}\,\Fmat$, or $\det\Fmat$ \\
Anisotropic & $\{\bar{I}_1\}$, $\{\bar{I}_2\}$, $\{J\}$, $\{I_4, I_5\}$
  & Fiber invariants are convex in $\Fmat$ \\
\bottomrule
\end{tabular}
\end{table}

\subsection{Architecture summary}

\begin{table}[htbp]
\centering
\caption{Summary of neural architectures. All use softplus activations and
  two hidden layers of width~64 unless noted.}
\label{tab:architectures}
\begin{tabular}{@{}lcccl@{}}
\toprule
\textbf{Architecture} & \textbf{Convex?} & \textbf{Polyconvex?}
  & \textbf{Parameters} & \textbf{Notes} \\
\midrule
MLP            & No  & No  &
  % TODO: Fill actual parameter counts from benchmarks.
  \textbf{[N]}
  & Universal baseline \\
ICNN           & Yes & No  &
  \textbf{[N]}
  & Convex in invariant space \\
Polyconvex     & Yes & Yes &
  \textbf{[N]}
  & Separate ICNN per invariant group \\
\bottomrule
\end{tabular}
\end{table}

% ── Architecture diagram placeholder ──────────────────────────────
\begin{figure}[htbp]
\centering
% TODO: Replace with TikZ or PDF diagram showing:
%   (a) MLP: x -> hidden -> hidden -> W
%   (b) ICNN: x -> hidden (non-neg z-weights, skip connections) -> W
%   (c) Polyconvex: x split into groups -> branch ICNNs -> sum -> W
\fbox{\parbox{0.9\textwidth}{\centering\vspace{4cm}
  \textbf{[PLACEHOLDER]} \\[0.5em]
  Architecture diagrams: (a) MLP, (b) ICNN with skip connections and
  non-negative z-path weights, (c) Polyconvex ICNN with separate branches
  per invariant group.
\vspace{1cm}}}
\caption{Schematic of the three neural architectures.  Dashed arrows in (b)
  indicate skip connections from the input; $\oplus$ symbols with ``$\geq 0$''
  labels indicate non-negative weight enforcement via softplus
  reparameterization.}
\label{fig:architectures}
\end{figure}
